% Generated draft: requires mathematical and source review.
\documentclass[11pt]{article}
\usepackage[margin=25mm]{geometry}
\usepackage{fontspec}
\setmainfont{Latin Modern Roman}
\usepackage{amsmath,amssymb,mathtools,unicode-math}
\setmathfont{Latin Modern Math}
\usepackage{enumitem,needspace,xurl,hyperref}
\setlength{\parindent}{0pt}
\setlength{\parskip}{5pt}
\setlength{\emergencystretch}{3em}
\newcommand{\N}{\mathbb{N}}
\newcommand{\Z}{\mathbb{Z}}
\newcommand{\Q}{\mathbb{Q}}
\newcommand{\R}{\mathbb{R}}
\newcommand{\C}{\mathbb{C}}
\newcommand{\F}{\mathbb{F}}
\newcommand{\A}{\mathbb{A}}
\newcommand{\PP}{\mathbb{P}}
\newcommand{\E}{\mathbb{E}}
\newcommand{\eps}{\varepsilon}
\newcommand{\dd}{\,\mathrm{d}}
\begin{document}
\section*{UnsolvedMath ID: 30003088}
\textbf{Source problem code:} catalogue entry
\subsection{Definitions and mathematical statement}
The supplied record does not define ``strictly complex'' or specify the construction of the logarithmic Chern numbers. The following gives a conditional formulation using the standard projective-plane interpretation; these conventions require confirmation against the source.

Let \(\mathcal L=\{L_1,\ldots,L_d\}\) be a finite arrangement of distinct lines in \(\PP^2_{\C}\). Blow up every intersection point incident with at least three lines, obtaining \(\pi:X\to\PP^2_{\C}\), and let \(D\) be the reduced total transform of \(\bigcup_i L_i\). Thus \(D\) has simple normal crossings. Define
\[
\bar c_1^2(\mathcal L)=(K_X+D)^2,
\qquad
\bar c_2(\mathcal L)=\deg c_2\bigl(\Omega_X^1(\log D)\bigr).
\]
For arrangements with \(\bar c_2(\mathcal L)\ne0\), define their logarithmic Chern slope by
\[
s(\mathcal L)=\frac{\bar c_1^2(\mathcal L)}{\bar c_2(\mathcal L)}.
\]
The source does not supply further admissibility hypotheses, such as positivity of \(\bar c_2\), exclusions of pencils, or bounds on intersection multiplicities; none is asserted here.

Let \(\mathscr T\) denote the source-intended class of strictly complex arrangements. This class remains unresolved: failure of a given complex arrangement to become defined over \(\R\) after a projective change of coordinates is different from failure of its incidence structure to admit any real realization. The record does not specify which condition is intended. Subject to resolving this qualification, put
\[
S_{\mathscr T}=\{s(\mathcal L):\mathcal L\in\mathscr T,\ \bar c_2(\mathcal L)\ne0\}.
\]
For \(S\subseteq\R\), define its set of accumulation points by
\[
\operatorname{Acc}(S)=\{a\in\R:\ \forall\varepsilon>0\ \exists s\in S\text{ with }0<|s-a|<\varepsilon\}.
\]
\textbf{Problem, with the preceding qualifications.} Determine
\[
\operatorname{Acc}(S_{\mathscr T})\cap\left(\frac52,\frac83\right],
\]
and determine how these accumulation phenomena depend on fields of definition of the arrangements. Here the lower endpoint \(5/2\) is excluded and the upper endpoint \(8/3\) is included.

For a subfield \(K\subseteq\C\), one possible precise meaning of being defined over \(K\) is that, after a projective change of coordinates over \(\C\), every line has a nonzero homogeneous linear equation with coefficients in \(K\). Under that convention one can form
\[
S_{\mathscr T}(K)=\{s(\mathcal L):\mathcal L\in\mathscr T,\ \bar c_2(\mathcal L)\ne0,\ \mathcal L\text{ is defined over }K\}
\]
and study \(\operatorname{Acc}(S_{\mathscr T}(K))\cap(5/2,8/3]\). This supplies a possible framework for the field question; the record specifies neither the fields to compare nor a conjectured relation. It also leaves unresolved whether slopes must themselves lie in \((5/2,8/3]\) before taking accumulation points, and whether limits obtained from distinct arrangements with an unchanged slope are intended to count.
\subsection{Short English statement}
Determine which values in \((5/2,8/3]\) are accumulation points of logarithmic Chern slopes of strictly complex line arrangements, and describe any dependence on their fields of definition. The meanings of ``strictly complex,'' the admissible arrangements, and the intended field comparison require confirmation. The supplied registry labels the problem open; this draft does not independently verify that status.
\subsection{Sources}
\begin{itemize}
\item Supplied local UnsolvedMath registry record, permanent ID \(30003088\), source code \texttt{catalogue entry}, title: \emph{Chern Slope Accumulation for Complex Line Arrangements}. The record belongs to the Oberwolfach Reports Open Problems collection and assigns the year 2016. No individual proposer is supplied.
\item Supplied report attribution: \emph{Mini-Workshop: Arrangements of Subvarieties, and their Applications in Algebraic Geometry} (2016), Oberwolfach Reports, pp.\ 653--702, DOI \texttt{10.4171/owr/2016/14}. \url{https://doi.org/10.4171/owr/2016/14}. The background describes these reports as documenting mathematical workshops and their open-problem discussions. No problem-specific page is supplied.
\item Supplied literature source, labeled Czapliński--Pokora: \url{https://arxiv.org/abs/1511.03879}. No title or specific supporting passage is supplied in the record.
\item Supplied literature source, labeled Urzúa survey: \url{https://www.scielo.cl/pdf/proy/v41n2/0716-0917-proy-41-02-517.pdf}. No specific supporting passage is supplied in the record.
\item Literature assessment dated exactly \texttt{2026-08-21}: the supplied curation labels the problem open and reports dense accumulation ranges up to \(5/2\) over the reals and finite-field analogues, while describing accumulation points in the strictly complex interval \((5/2,8/3]\) as an open problem. This is supplied dated curation derived from the OWR record, not independently verified literature assessment. The record explicitly calls for independent verification of time-sensitive status and cited claims.
\item The registry labels the statement \texttt{corrected\_verified}; its accompanying explanation limits the reported verification to repairs of notation, typography, and sentence structure without changing mathematical content. This does not establish mathematical review of the conventions used in this draft.
\end{itemize}
\end{document}
